\section{Conclusion}

Tool-call interfaces censor agent trajectories, silently and per checkpoint; within Qwen2.5-Coder
the absolute undercount rises monotonically with checkpoint size, which we report as a
within-family observation rather than a general law. In a 150-step RL control differing only in
verl's parser registry, the mismatched arm turns 8{,}689 tight-classified call-like emissions into 10 executions,
whereas the repaired arm executes and observes 16{,}844 calls. The intervention therefore restores
the tool-mediated experience distribution. We detect no held-out improvement: rescues remain 12
and final passes move from 430/542 to 427/542. Access to a trajectory and learning from it are
separate empirical questions.

Two controls fix what the effect is a property of. A $2\times2$ over chat template and parser on a
public benchmark finds zero recovery from either one-component replacement at the documented
baseline and 196/200 parses under joint replacement; and a supporting ladder under a
\emph{matched} interface holds the silent fraction at 0--2 across
a comparable span of scale, where the mismatched interface takes it to 80. \textbf{What the serving
layer removes is not a property of the model, nor of the parser, but of the contract between them.}

The most uncomfortable finding is methodological. We ran three single-variable paired controls on a
23\% failure and concluded it was a model deficiency; a fourth, taken from a sentence in the serving
documentation, reduced it to zero. Our own human validation reproduced the same mechanism: twelve of
ninety-eight annotations missed a tool call because it sat after a fenced code block, and the
annotator --- like the parser --- concluded from the same bytes that no call had been made.
Reporting server-parsed tool-call rates as model capability is not a safe default.
